\chapter{Discussion : que signifie la rareté du langage doctrinal ?}\label{ch:rarete}

\section{Le constat}

Les résultats des chapitres précédents peuvent être résumés ainsi. Sur \st{nC} textes funéraires, \st{nProfilLarge} désignent Dieu, le Christ, l’Esprit ou la Trinité par un mot lu ou abrégé ; une règle plus restrictive en retient \st{nProfilRestr}, la règle stricte \st{nProfilStrict}. Le Christ est nommé en toutes lettres dans un seul texte. L’Esprit saint n’est nommé, au plus, qu’une fois, sur restitution partielle ; la Trinité une fois, sans énumération des personnes ; la résurrection jamais de manière assurée. Le vocabulaire le plus fréquent est celui du repos et de la paix ; viennent ensuite les qualifications de la personne (\emph{fidelis}, \emph{neophytus}, \emph{ancilla dei}, \emph{virgo}) et, dans quelques poèmes, les images de la lumière, du ciel et de la société des saints. Le chrisme et la croix, en revanche, sont décrits par l’EDB sur \st{nSignesC} textes (chapitre~\ref{ch:supports}).

Ces comptes décrivent des catégories de lecture, non un taux d’adhésion doctrinale. Ils sont robustes au traitement des restitutions : la part des textes concernés varie entre 1,7 et 3,9~\% selon la règle (chapitre~\ref{ch:typologie}).

\section{Trois lectures possibles}

Ce constat admet au moins trois interprétations, que le corpus ne permet pas de départager.

\emph{Une lecture générique.} L’épitaphe est un genre contraint, qui privilégie l’identification du défunt, la datation et l’éloge. Sa brièveté laisse peu de place à la profession de foi ; son destinataire, la communauté des vivants qui visitent la tombe, partage les convictions qu’elle n’a pas besoin d’énoncer. Le monogramme, la formule \emph{in pace} et l’emplacement dans un cimetière chrétien disent l’appartenance ; la doctrine n’a pas à y être exposée. Dans cette perspective, les rares textes doctrinaux sont des exceptions dues à des commanditaires ou à des rédacteurs particuliers : l’évêque Damase pour Proiecta (\nnum{244}), le milieu sénatorial de Lucceia (\nnum{306}), l’entourage ascétique de Quintilianus (\nnum{474}).

\emph{Une lecture sociale.} Les textes les plus riches en vocabulaire religieux appartiennent souvent à des milieux aisés (sarcophages, poèmes, rang sénatorial) ou cléricaux. La rareté du langage doctrinal pourrait refléter une différence de culture religieuse et de culture écrite entre ces milieux et la majorité des commanditaires. Mais le corpus documente trop peu le statut des commanditaires pour vérifier cette hypothèse, et plusieurs textes simples (\emph{ancilla dei}, \emph{fidelis}, \emph{neophytus}) expriment une identité religieuse sans éloquence.

\emph{Une lecture documentaire.} Une part de ce que nous lisons est l’œuvre des éditeurs. Les restitutions concernent en priorité les formules attendues, et les écarts entre les deux bases (\nnums{73} et 447 en particulier) montrent que l’énoncé religieux peut dépendre d’un choix de restitution. La sélection des textes datés, enfin, n’est pas aléatoire : elle favorise les textes complets, donc souvent les plus longs. Ces biais jouent en sens contraires ; ils n’expliquent pas à eux seuls la rareté observée.

\section{Ce que l’on ne peut pas conclure}

Il serait anachronique de déduire de l’absence de formules trinitaires une absence de foi trinitaire, comme il le serait de lire dans \emph{in pace Christi} une christologie. Le silence d’un genre n’est pas le silence d’une communauté. Les décennies 340–403 sont, à Rome, celles des controverses ariennes, du schisme d’Ursinus et de Damase, du concile de Constantinople et de la législation théodosienne ; aucun de ces événements n’est perceptible dans les épitaphes datées, sinon peut-être, à la toute fin de la période, dans le \emph{confirmans trinitatem} du \nnum{474}.

Le corpus permet de décrire que, dans les épitaphes romaines datées de 340 à 403, la foi s’exprime surtout par des formules de repos, de paix et d’appartenance, par des signes, et très rarement par des énoncés doctrinaux. La comparaison avec la littérature patristique contemporaine, qui donnerait la mesure de cet écart, et avec les épitaphes non datées des mêmes cimetières, qui permettrait de vérifier si les textes datés sont représentatifs, sort du cadre de ce volume. Elle supposerait des corpus de référence définis avec la même rigueur.
